\documentclass[12pt]{article}
\usepackage{amsfonts,amsthm,amsmath,amssymb,mathtools}
\usepackage{mathrsfs}
\usepackage{enumitem}
\usepackage{tikz-cd}
\usepackage{geometry}
\usepackage{mlmodern}
\usepackage{minted}

\geometry{letterpaper, portrait, margin=1in}

\setcounter{section}{0}

\renewcommand{\thesection}{\S\arabic{section}}
\renewcommand{\thesubsection}{\arabic{section}}
\newcommand{\dd}{\mathop{}\!\mathrm{d}}
\newcommand{\R}{\mathbb{R}}
\newcommand{\Z}{\mathbb{Z}}
\newcommand{\C}{\mathbb{C}}
\newcommand{\Q}{\mathbb{Q}}

\DeclareMathOperator{\lcm}{lcm}
\DeclareMathOperator{\im}{im}

\newtheorem{thm}{Theorem}
\newenvironment{theorem}[1]{
	\renewcommand{\thethm}{#1}
	\begin{thm}
		}{\end{thm}}

\newtheorem{lma}{Lemma}
\newenvironment{lemma}[1]{
	\renewcommand{\thelma}{#1}
	\begin{lma}
		}{\end{lma}}

\theoremstyle{definition}
\newtheorem{ex}{}
\newenvironment{exercise}[1]{
	\renewcommand{\theex}{#1}
	\begin{ex}
		}{\end{ex}}

\title{Project 2: Edwards curves}
\author{Connor Mooneyhan}
\date{December 9, 2025}

\begin{document}

\maketitle

Prompt: Research Edwards curves and twisted Edwards curves and doing arithmetic on them over finite fields.
This is an alternative to a Weierstrass equation.
What advantages do these curves have in terms of speed of computation and in terms of use in the elliptic curve Diffie-Hellman key exchange mechanism? Learn about side-channel attacks on cryptosystems and how arithmetic with Edwards curves can reduce the potential of such attacks.
Write your own code that will implement arithmetic on Edwards curves using Magma (or Sage or Python or whatever coding platform you prefer).
I recommend the paper “Faster Addition and Doubling on elliptic curves” by Daniel Bernstein and Tanja Lange.
You should have the background you need for this by the end of Week 12.

\section{Introduction}
While the Weierstrass form of a curve is widely used, there are many other forms into which elliptic curves can be transformed through bi-rational equivalence.
Some of these may have advantages in terms of how easy it is to define and compute with their group law, as well as their consistency and regularity.
Edwards curves were discovered in 2007 by Harold Edwards, and they offer many of these advantages, as I have outlined below.

\section{Edwards Form}
Edwards showed that, similar to other forms like Weierstrass form, any elliptic curve over a finite field with more than two elements can be transformed into an equation of the form \begin{equation*}
	x^2 + y^2 = 1 + dx^2y^2
\end{equation*}
where $d$ is a point in the base field which is neither 0 nor 1.
The addition law is very simple, compared to the usual addition law for elliptic curves: \begin{equation*}
	(x_1, y_1) + (x_2, y_2) = \left( \frac{x_1y_2 + y_1x_2}{1+dx_1x_2y_1y_2}, \frac{y_1y_2-x_1x_2}{1-dx_1x_2y_1y_2} \right).
\end{equation*}
The identity is $(0, 1)$.

One can generalize Edwards curves to what are called twisted Edwards curves.
Let $k$ be the base field with order $> 2$, and let $a, d \in k$.
Then the twisted Edwards curve with coefficients $a$ and $d$ is the curve \begin{equation*}
	E_{E,a,d} : ax^2 + y^2 = 1+dx^2y^2.
\end{equation*}
Thus an Edwards curve is just a twisted Edwards curve with $a=1$.

\section{Advantages}
Given the explicit group law written in terms of just the coordinates of the addends and the constant $c$, it is certainly easy to write code that will perform addition on any given Edwards curve.
Not only that, but these curves are also much faster to compute.
In fact, as of the Bernstein and Lange paper from 2007, addition on Edwards curves was in contention for the fastest known method for performing certain calculations.

One use of it is in the curve known as Curve25519, which is bi-rationally equivalent to the Edwards curve \begin{equation*}
	x^2 + y^2 = 1 + \frac{121665}{121666}x^2y^2.
\end{equation*}
This is an officially approved elliptic curve for use with Diffie-Hellman key exchanges, and it is known for its speed in certain applications.

Another application is in protecting against side-channel attacks.
These are attacks that take advantage of side-effects of computation such as the time it takes to do computations, the noise a machine gives off, the heat it produces, etc.
Edwards curve computations are resistant to many of the most prevalent side-channel attacks due to their uniformity and regularity.

\section{Implementation}
We conclude with an implementation of a few different algorithms created by Bernstein and Lange, written in JavaScript.
The choices for $c$ and $d$ are of course arbitrary.
\begin{minted}{javascript}
const c = 1;
const d = 2;

// add two points
function add([X1, Y1, Z1], [X2, Y2, Z2]) {
  A = Z1 * Z2;
  B = A ** 2;
  C = X1 * X2;
  D = Y1 * Y2;
  E = d * C * D;
  F = B - E;
  G = B + E;

  const X3 = A * F * ((X1 + Y1) * (X2 + Y2) - C - D);
  const Y3 = A * G * (D - C);
  const Z3 = c * F * G;

  return [X3, Y3, Z3];
}

// double a point
function double(P) {
  return add(P, P);
}

// return a y-value for a given x-value if one exists
function compress(x) {
  const y = Math.sqrt((c ** 2 - x ** 2) / (1 - c ** 2 * d * x ** 2));

  return [x, y];
}
\end{minted}

\section{Bibliography}
\textit{Faster addition and doubling on elliptic curves} by David Bernstein and Tanja Lange
\\
Used for learning about the Edwards form of an elliptic curve and its benefits.
\\\\
\textit{Twisted Edwards Curves} by David Bernstein et. al.
\\
Used for learning about twisted Edwards curves.
\\\\
\textit{Edwards curves} by Youssef El Housni
\\
Used for understanding the application to Diffie-Hellman.
\\\\
https://www.wired.com/story/what-is-side-channel-attack/
\\
Used for learning about side-channel attacks.
\\\\
https://www.sciencedirect.com/science/article/abs/pii/S0141933118300516
\\
Used for learning about the benefits of Edwards curves in protecting against side-channel attacks.

\end{document}
